% !TEX root = handbook-physics.tex
% !TEX TS-program = pdflatex
% !TEX encoding = UTF-8 Unicode






\chapter{Elastic Waves}
\label{ch-30}

\section{Longitudinal and Transverse Waves in an Elastic Medium}
\label{sec-30.1} 

\subsection{}\label{30.1.1} A body is said to be \redem{elastic} and its deformation, due to the application of external forces, is \redem{elastic deformation} if it completely disappears after the external forces are removed. According to \redem{Hooke's law} elastic deformation is directly proportional to the external effects that cause it, i.e. it has a linear dependence on these effects. At sufficiently small deformation, all bodies can he assumed, practically, to be elastic.

The elastic properties of bodies depend upon the nature of the thermal
motion of their molecules and their forces of interaction. For
example, a gaseous body changes its shape without hindrance to suit
the shape of the vessel it is contained in; gas has no \redem{elasticity of
shape}. But a gas has \redem{elasticity of volume}, i.e. the capacity to resist
changes in its volume. This property of a gas is due to the thermal
motion of its molecules and is manifested in the change in the
pressure $p$ of the gas when its volume $V$ is changed.

According to \redem{Hooke's law for volume deformation}, the change
$\dd p$ in the pressure of a gas for a small change $\dd V$ in its
volume is directly proportional to the relative volume defor­ mation,
called the volume strain. Thus
\begin{equation*}%
\dd p = -K \frac{\dd V}{V},
\end{equation*}
where $K$ is called the \redem{bulk modulus of elasticity} of the gas. For an ideal gas \ssect{8.4.1} the value of $K$ depends upon the kind of thermodynamic process of compression \ssect{or expansion} being carried out. At a very slow change in the volume of the gas, the process can be considered isothermal \ssect{8.3.7}, whereas at a rapid change it can be considered adiabatic \ssect{8.3.7}. In the former case $K= p$ and in the latter $K = \varkappa p$, where $\varkappa$ is the adiabatic exponent \ssect{9.5.11}.

\subsection{}\label{30.1.2} The elasticity of a crystalline solid is due to the forces of mutual attraction and repulsion of its particles (ions, atoms or molecules) that make up the body and that have random thermal vibration about the points of its crystal lattice. The forces of interaction of the particle impede deformation of the crystal lattice that is associated with both the volume and shape of the body. Therefore, in addition to their bulk elasticity, solids also have elasticity of form, which is manifested by their resistance to shear deformation \ssect{40.3.9}.

The elasticity of liquids is also due to the forces of intermolecular interaction. But, owing to the fact that the average time $\tau$ that molecules of a liquid remain in a temporarily settled state \ssect{13.2.4} is very short, liquids, like gases, possess only bulk elasticity. They manifest elasticity of form only with respect to variable deformation of ultrahigh frequency with a period that is less than or of the order of $\tau$.

\subsection{}\label{30.1.3} \redem{Elastic waves} are \redem{mechanical disturbances} (deformations) propagated in an elastic medium. Bodies that act on the medium and cause these disturbances are called sources of the waves. A member of
the audience in a theatre, for instance hears the speech and singing of the actors and the sounds of the musical instruments due to the variations in the pressure of the air, reaching the audience and caused by the sources of sound. 

\redem{Sonic}, or \redem{acoustic, waves} are elastic waves of low intensity, i.e. weak mechanical disturbances propagated in an elastic medium.

Sonic waves, acting on the organs of hearing of people, are capable of producing auditory sensations if the frequency $\nu$ of their
corresponding vibrations lies in the range from 16 to \SI{2d4}{\hertz} (\redem{audible sound}). Elastic waves with a frequency $nu < \SI{16}{\hertz}$ are said to be \redem{infrasonic}, whereas those with a frequency $\nu > \SI{2d4}{\hertz}$ are \redem{ultrasonic} (elastic waves) with $\nu \gg \SI{d9}{\hertz}$ are often said to be \redem{hypersonic}.

\subsection{}\label{30.1.4} The propagation of elastic waves in a
medium does not involve the transport of matter.\sidenote{A certain transport of matter (mass transfer) may occur in the propagation in a medium of strong disturbances (for example, shock waves in explosions), when the vibration of the particles of the medium becomes nonlinear.} In a boundless medium, it consists in exciting forced vibrations of parts of the medium that are farther and farther away from the source of the waves. It is feasible here to digress from the discrete (molecular) structure of the medium, dealing with it as a \redem{continuous medium}, distributed continuously in space and possessing definite elastic properties. 

A \redem{particle} of such a medium, having forced vibration, is understood to be a small element of its volume whose size, however, is many times greater than the intermolecular distances so that it contains an enormous number of molecules. Practically, such a particle can be dealt with as a point because even in a gas the intermolecular distances are extremely small (of the order of \SI{d-8}{\meter} under standard conditions).

\subsection{}\label{30.1.5} An elastic wave is said to be \redem{longitudinal} if the particles of the medium vibrate in the direction of wave propagation. Longitudinal waves are associated with bulk deformation of the elastic medium and can therefore propagate in any kind of medium: solid, liquid or gaseous. An example of such waves are the sonic waves in air.

An elastic wave is said to be \redem{transverse} if the particles of the medium vibrate, remaining in planes perpendicular to the direction of wave propagation. Transverse waves are associated with shear deformation of the elastic medium, and, consequently, can be generated and propagated only in media having elasticity of shape, i.e. in solids. An example of transverse waves are
those propagated along the strings of musical instruments. 

Special properties are found in \redem{surface} waves, which propagate along the free surfaces of liquids (or the interface between two immiscible liquids), being a perturbance of the surface due to external effects (dropping of bodies into the liquid, the motion of ships, the wind, etc.). The forces of gravity and surface tension play a decisive role in the generation and propagation of these waves. In surface waves the particles of liquid have transverse and longitudinal vibrations simultaneously, describing elliptical or more complex paths.

\subsection{}\label{30.1.6} A medium is said to be \redem{uniform} if its physical properties, essential in the problem being considered, do not change from point to point. A medium that is uniform with respect to certain physical properties may be non-uniform with respect to others. A mono-crystalline body, for example, is uniform with respect to its elastic properties and, at the same time, optically non-uniform for X-rays.

A medium is \redem{isotropic} if its physical properties that are essential in the problem being considered are the same in all directions. A medium that is isotropic with respect to certain physical properties may be anisotropic with respect to others. Cubic crystals, for instance, are optically isotropic, but are anisotropic with respect to other properties. Gases and liquids are isotropic with respect to any physical properties.

\subsection{}\label{30.1.7} A medium is said to be \redem{linear} if there is a directly proportional relation between the quantities characterizing the external effects being considered that act on the medium and the changes they lead to in the state of the medium. An elastic medium, for example, obeying Hooke’s law \ssect{30.1.1} is linear with respect to its mechanical properties. A dielectric is a linear medium with respect to its electrical properties if its dielectric constant \ssect{14.2.4} is independent of the electric field strength. By analogy, a magnetic material is a linear medium with respect to its magnetic properties if its relative magnetic permeability \ssect{26.4.5} is independent of the magnetic induction of the field.